\chapter{Metodologia}
\label{cap:metodologia}

\notaguia{Primeira redação, a partir da Parte~4 da visualização (``A
oficina''), da Parte~3 do \texttt{guia\_de\_leitura.md} e dos documentos de
\texttt{Textos/metodologia/}. As referências de método antes citadas apenas
por epônimo foram conferidas e vinculadas, inclusive o artigo original de
Weiszfeld (1937) e a citação formal da Coleção~10.1 do MapBiomas, que eram as
duas pendências desta seção.}

Este capítulo declara o desenho da pesquisa: área de estudo, fontes, unidades
de análise, painel, régua temporal, instrumental de estimação e testes de
robustez. Duas seções fogem ao padrão de um capítulo de método e existem por
razões específicas deste trabalho: a Seção~\ref{sec:met-decisoes} registra as
decisões que alteraram resultados já obtidos, com o veredito anterior e o
motivo da mudança, e a Seção~\ref{sec:met-ia} declara o papel da inteligência
artificial como instrumento de trabalho.

O trabalho é observacional e quantitativo, apoiado inteiramente em dados
públicos. Não há levantamento primário: a contribuição está no cruzamento
sistemático de bases já existentes, na construção de unidades espaciais e
temporais comparáveis e na disciplina de teste aplicada a cada afirmação.

\section{Área de estudo}
\label{sec:met-area}

A área de estudo é o estado de Goiás, integralmente inserido no bioma
Cerrado, em seus 246 municípios atuais. O recorte estadual não é apenas
administrativo: é ele que dá acesso simultâneo à estatística socioeconômica
municipal --- publicada pelo IBGE e pelo Banco Central nessa malha --- e à
série de sensoriamento remoto de resolução fina.

Goiás atravessa, ainda, um gradiente latitudinal que vai de um Sul agrícola
consolidado a um Norte de ocupação mais recente, com o remanescente de
Cerrado concentrado nessa porção setentrional. É essa heterogeneidade interna
que permite testar, dentro de uma mesma unidade federativa, os dois regimes
de fronteira que a literatura do bioma descreve em escala regional
(Seção~\ref{sec:ref-matopiba}).

O eixo Sul--Norte é operacionalizado pelas cinco mesorregiões do IBGE (malha
de 2017, decisão D6), atribuídas a cada município por junção espacial sobre
os centroides municipais. Elas são ordenadas pela latitude média dos pixels
de conversão, e é essa ordenação --- não a nomenclatura --- que define o eixo
analítico (Tabela~\ref{tab:mesorregioes}).

\begin{table}[htb]
\centering
\caption[Mesorregiões de Goiás, ordenadas por latitude]{Mesorregiões de Goiás (IBGE, 2017), ordenadas por latitude média
dos pixels de conversão.}
\label{tab:mesorregioes}
\small
\begin{tabular}{lrr}
\toprule
\textbf{Mesorregião} & \textbf{Municípios} & \textbf{Latitude média} \\
\midrule
Sul Goiano      & 82 & $-17{,}7^\circ$ \\
Centro Goiano   & 82 & $-16{,}0^\circ$ \\
Leste Goiano    & 32 & $-15{,}7^\circ$ \\
Noroeste Goiano & 23 & $-15{,}2^\circ$ \\
Norte Goiano    & 27 & $-14{,}1^\circ$ \\
\bottomrule
\end{tabular}
\fontefig{elaboração própria a partir da malha de mesorregiões do IBGE (2017).}
\end{table}

A alternativa de usar as Regiões Geográficas Imediatas (133 unidades em
Goiás) foi descartada: pulverizaria o painel sem ganho de leitura no eixo
Sul--Norte. A Figura~\ref{fig:localizacao} situa o estado e a partição
mesorregional adotada.

Uma segunda partição, mais grossa, aparece em parte dos resultados, e
convém defini-la aqui para que os rótulos não se confundam. Quando o texto
falar em \textbf{três regiões} --- Sul, Centro e Norte ---, trata-se do
agrupamento das cinco mesorregiões em três blocos latitudinais: o Sul
corresponde ao Sul Goiano; o Centro reúne o Centro Goiano e o Leste Goiano;
e o Norte reúne o Norte Goiano e o Noroeste Goiano. Essa malha de três é
usada onde o número de unidades por bloco precisa ser grande o bastante para
sustentar comparação --- a decomposição da fronteira, o desenho de
\emph{drive} comum e a conta de carbono. Onde o texto disser
``mesorregião'', a malha é a de cinco. Como ``Sul'' nomeia a mesma coisa nas
duas partições, mas ``Centro'' e ``Norte'' não, cada resultado adiante
declara a malha em que foi medido.

\begin{figure}[htb]
\centering
\caption[Área de estudo e divisão em mesorregiões]{Área de estudo: Goiás no Brasil e no bioma Cerrado, e a divisão em
mesorregiões sobre a malha municipal.}
\label{fig:localizacao}
\includegraphics[width=\textwidth]{cap3_localizacao.pdf}
\fontefig{elaboração própria a partir das malhas do IBGE obtidas pelo pacote
\texttt{geobr}: mesorregiões de 2017 --- a mesma malha que a análise emprega
(decisão D6) ---, municípios e estados de 2020, biomas de 2019. Projeção
Policônica do Brasil (SIRGAS 2000, EPSG:5880); os dois painéis estão em
escalas diferentes, cada um com a sua régua. O preenchimento das
mesorregiões segue a ordenação por latitude média dos pixels de conversão da
Tabela~\ref{tab:mesorregioes} --- do Sul (mais claro) ao Norte (mais
escuro) ---, e não a nomenclatura. O vazio a leste do Centro Goiano é o
Distrito Federal, que não integra o estado.}
\end{figure}

\section{Fontes de dados}
\label{sec:met-fontes}

Todas as fontes são públicas e foram coletadas por rotinas automatizadas,
com cache local dos arquivos brutos, de modo que qualquer número do trabalho
possa ser reconstruído a partir do dado original. A
Quadro~\ref{quadro:fontes} resume fonte, cobertura e papel na análise.

A espinha dorsal é a série anual de uso e cobertura da terra do
\citeonline{MapBiomas2025}, em sua Coleção 10.1. A ela somam-se os dados de
pesquisa e registro administrativo do \citeonline{IBGE_SIDRA}, o crédito
rural do \citeonline{BACEN_SICOR}, o índice de desenvolvimento municipal da
\citeonline{FIRJAN_IFDM}, o mapa de aptidão agrícola da
\citeonline{Embrapa_Aptidao}, o cadastro de unidades armazenadoras da
\citeonline{CONAB_Armazenagem}, o monitoramento de desmatamento do
\citeonline{INPE_PRODES} --- usado como aferição externa, e não como insumo
---, o Atlas do Desenvolvimento Humano \cite{IPEA_IDHM} e as malhas
territoriais oficiais obtidas pelo pacote \emph{geobr}
\cite{Pereira2019geobr}.

As siglas nela empregadas designam, pela ordem: o Sistema IBGE de Recuperação
Automática (SIDRA), portal pelo qual se acessam as pesquisas do Instituto
Brasileiro de Geografia e Estatística (IBGE) --- entre elas a Produção
Agrícola Municipal (PAM), a Pesquisa da Pecuária Municipal (PPM), as Contas
Regionais, de onde vêm o Produto Interno Bruto (PIB) e o Valor Adicionado
municipais, o Índice Nacional de Preços ao Consumidor Amplo (IPCA) e o Censo
Agropecuário; o Sistema de Operações do Crédito Rural e do Proagro (SICOR),
do Banco Central do Brasil (BACEN); o Índice de Desenvolvimento Humano
Municipal (IDH-M), publicado pelo Instituto de Pesquisa Econômica Aplicada
(IPEA); a plataforma Trase, iniciativa de rastreamento de cadeias de
\emph{commodities}; a Empresa Brasileira de Pesquisa Agropecuária (Embrapa);
e o PRODES, sistema de monitoramento do desmatamento por satélite do
Instituto Nacional de Pesquisas Espaciais (INPE).

\begin{quadro}[htb]
\centering
\caption[Fontes de dados: cobertura e papel na análise]{Fontes de dados: cobertura temporal e papel na análise.}
\label{quadro:fontes}
\small
\begin{tabular}{p{3.6cm}p{2.4cm}p{7.6cm}}
\toprule
\textbf{Fonte} & \textbf{Cobertura} & \textbf{Papel na análise} \\
\midrule
MapBiomas, Coleção 10.1 (raster, 30~m) & 1985--2024, anual &
Uso e cobertura da terra: estoques por classe, matrizes de transição
pixel a pixel, centros de massa, idade da pastagem. É a espinha dorsal
do trabalho \\
\addlinespace
IBGE/SIDRA --- PAM (tabelas 1612, 1613 e 839) & 1985--2024 &
Área plantada, área colhida, produção e rendimento das lavouras; milho de
1ª e 2ª safras (2003--2024) \\
\addlinespace
IBGE/SIDRA --- PPM (tabelas 3939, 74 e 94) & 1985--2024 &
Efetivo de rebanhos (bovinos e demais), leite e ovos \\
\addlinespace
IBGE/SIDRA --- Contas Regionais (tabela 5938) & 1999--2021 &
PIB municipal e Valor Adicionado por setor \\
\addlinespace
IBGE/SIDRA --- tabelas 6579 e 1737 & 2001--2024; mensal &
População residente; IPCA (deflator) \\
\addlinespace
IBGE/SIDRA --- Censo Agropecuário 2017 & 2017 (transversal) &
Plantio direto, uso de calcário, orientação técnica --- variáveis de
manejo indisponíveis em série \\
\addlinespace
SICOR/BACEN & 2013--2026 &
Crédito rural por município, finalidade e produto \\
\addlinespace
IPEA --- IDH-M & 1991, 2000, 2010 &
Desenvolvimento humano municipal; série descontinuada após 2010 \\
\addlinespace
Trase & 2004--2020 &
Cadeia produtiva da soja e da carne bovina, com separação explícita entre
volume exportado e produção \\
\addlinespace
Embrapa --- aptidão agrícola & estática &
Exposição física exógena (gradiente edafoclimático) nos desenhos de
interação \\
\addlinespace
PRODES/INPE & 2001--2024 &
Validação externa das séries de supressão de vegetação \\
\bottomrule
\end{tabular}
\fontefig{elaboração própria.}
\end{quadro}

Duas observações sobre as variáveis monetárias. Todo valor em reais --- PIB,
Valor Adicionado, crédito --- é deflacionado para reais de dezembro de 2024
pelo IPCA acumulado, construído como número-índice a partir da série mensal e
avaliado em dezembro de cada ano. E as variáveis monetárias anteriores ao
Plano Real foram excluídas da coleta em vez de convertidas: a conversão entre
Cruzeiro, Cruzado, Cruzado Novo e Cruzeiro Real introduziria erro de medida em
séries usadas como regressores. Por isso o valor da produção da PAM e a série
de leite em valor não entram no painel, embora as quantidades físicas
correspondentes entrem.

\section{Unidades de análise: municípios e Áreas Mínimas Comparáveis}
\label{sec:met-amc}

A malha municipal brasileira não é constante no tempo, e Goiás é um caso
severo: o estado tinha 184 municípios com estatística municipal em 1985 e tem
246 hoje. Dados
de sensoriamento remoto e dados de pesquisa comportam-se de maneiras opostas
diante de uma emancipação: o raster é recortado pelo polígono \emph{atual} em
todos os anos, ao passo que a estatística do IBGE é tabulada pelo município
\emph{como ele existia naquele ano}. Ao juntar as duas fontes pelo código
municipal, compara-se, para um município-pai, a cobertura da terra do
território de hoje contra a produção de um território historicamente maior;
para os municípios-filhos, compara-se cobertura contra dado ausente.

O efeito é grande. Dos 246 municípios, 62 (25\%) só aparecem na estatística
municipal depois de 1985, em ondas de emancipação concentradas em 1989 (27),
1993 (21), 1997 (10) e 2001 (4). O total estadual do rebanho atravessa esses
anos sem descontinuidade --- o gado do filho estava contabilizado no pai ---
mas municípios-pai individuais registram quedas de 50\% a 80\% no ano da
emancipação. Ler essas quedas como dinâmica pecuária seria atribuir a um
fenômeno econômico o que é perda de território. E como as ondas de 1993 e
1997 caem dentro das janelas usadas pelas análises longitudinais, tais quedas
dominariam qualquer estimador de variação nesse período.

A solução adotada (decisão D11) é a agregação em Áreas Mínimas Comparáveis
\cite{Ehrl2017}: cada município-pai é reunido aos seus desmembramentos numa
unidade de território constante ao longo de toda a janela. Aplicada a Goiás
com a concordância de 1980--2010 --- janela que antecede todas as
emancipações do período analisado ---, a agregação reduz os 246 municípios a
\textbf{166 AMCs}, sendo 53 grupos de pai e filhos e 113 unidades que
permanecem idênticas ao município. A validação é direta: o total estadual é
idêntico ao da malha municipal e nenhuma das 166 AMCs estreia depois de
1985.

A agregação obedece a uma regra que vale para qualquer variável derivada:
grandezas \emph{extensivas} (hectares, cabeças, reais) são somadas, e
grandezas \emph{derivadas} (razões, densidades, lotação) são
\emph{recalculadas} a partir das extensivas já somadas --- nunca somadas ou
promediadas diretamente, o que produziria média de razões em lugar de razão
de totais.

Disso resulta uma estratégia de dois trilhos, e usar o trilho errado é erro
metodológico: as AMCs (166 unidades) são a malha de toda análise
longitudinal --- primeiras diferenças, painel com efeitos fixos,
periodização, tendências, deslocamento de centros de massa ---, enquanto a
malha municipal atual (246 unidades) fica reservada às análises transversais
e ao período recente, como o Censo Agropecuário de 2017 e os mapas de um
único ano. O preço, com a aritmética correta: os 133 municípios reunidos nos
53 grupos deixam de ser analisáveis isoladamente nas séries longas, e a malha
perde 80 unidades em relação às 246. É um preço inevitável quando o dado vem
de pesquisa, porque a informação histórica do município-filho não existe
separada --- ela estava no pai.

Duas convenções geodésicas sustentam as medidas, e vale separá-las porque
respondem a perguntas diferentes. Toda operação \emph{geométrica} ---
centroides, distâncias, deslocamentos, recortes zonais --- é feita no sistema
de coordenadas de código EPSG:5880, Sistema de Referência Geocêntrico para as
Américas (SIRGAS 2000) em \textbf{projeção Policônica do Brasil}, que é o
padrão cartográfico brasileiro e a projeção em que saem os mapas deste texto;
o EPSG:4674 aparece apenas quando a leitura pede coordenadas geográficas. A
policônica é projeção de compromisso --- não é conforme nem equivalente ---,
e por isso ela não é usada para medir área. As áreas vêm de outro caminho: o
censo de pixels conta células de 30~m na grade nativa do produto e converte a
contagem em hectares pelo cosseno da latitude média observada de cada
município, sem passar por projeção nenhuma. A malha municipal de referência é
a do IBGE para 2020.

\section{O painel unificado e as primeiras diferenças}
\label{sec:met-painel}

As fontes descritas convergem para uma tabela única, organizada por
município e ano --- $246 \times 40 = 9.840$ linhas e 195 colunas --- que
reúne, para cada observação, os estoques de uso da terra, a produção
agrícola e pecuária, o produto e o valor adicionado, o crédito rural, a
população e as variáveis derivadas. É desse painel, e de sua versão agregada
em AMCs, que sai toda a estimação do trabalho.

A segunda decisão estruturante (D7) é que \textbf{nenhuma associação é
estimada em nível}: todas as variáveis entram em primeira diferença,
$\Delta x_t = x_t - x_{t-1}$. A razão é que quase toda série do painel tem
tendência --- área, produção, produto e população crescem ao longo de quatro
décadas --- e correlacionar níveis entre séries com tendência comum produz
associações altas e vazias. Diferenciar remove a tendência e reformula a
pergunta: em vez de ``municípios com mais crédito têm menos pasto?'', o
modelo responde ``quando o crédito de um município subiu, o pasto dele
recuou?''.

Essa escolha foi objeto de auditoria específica. Muitas variáveis do painel
são internamente relacionadas por construção: área plantada e produção, ou
área e valor adicionado, medem facetas do mesmo fato e apresentam correlação
em nível da ordem de 0,9. Estimar uma sobre a outra em nível não seria uma
associação --- seria uma identidade contábil disfarçada. A primeira
diferença, combinada com os efeitos fixos descritos adiante, desarma esse
vazamento: as mesmas duplas apresentam correlação próxima de zero em variação
\emph{within}. A auditoria de circularidade residual encontrou um único caso
não resolvido pela diferenciação --- o abate bovino, que na fonte é
reconstruído como taxa aplicada ao rebanho e, portanto, é redundante com ele
por definição --- e essa variável foi descartada como evidência. Onde ela
ainda aparece em resultado de pipeline, o coeficiente é declarado identidade
contábil, e não associação estimada.

Duas observações da série perdem-se na montagem do painel de estimação, e
convém dizer por quê antes que os $n$ divirjam aos olhos do leitor: a série
de cobertura tem 40 observações anuais (1985--2024), a primeira diferença
consome a primeira e a defasagem de um ano consome a segunda. Daí as
\textbf{38 diferenças anuais} que aparecem nas regressões, no cálculo de
erros-padrão robustos a autocorrelação e na contagem de realizações
independentes do \emph{shifter} macroeconômico. Quando o texto disser 40, o
objeto é a série descritiva; quando disser 38, é o painel estimado.

\section{Periodização orientada pelos dados}
\label{sec:met-periodizacao}

Narrar quarenta anos exige períodos, e a escolha das fronteiras entre eles é
decisão metodológica, não detalhe de exposição. O trabalho recusou ancorar os
períodos em marcos institucionais escolhidos pelo pesquisador --- a pergunta
``por que 2005 e não 2006?'' não teria resposta formal --- e inverteu a
lógica: as fronteiras são estabelecidas pelos dados, e os marcos
institucionais são classificados \emph{a posteriori} pela força da evidência
que os sustenta.

A regra de decisão foi fixada antes de rodar os testes: um método primário
--- estatística sup-F de Quandt-Andrews \cite{Quandt1960, Andrews1993}, na versão multivariada, aplicada
conjuntamente às variações anuais de vegetação natural, pastagem e
agricultura, com segmentação binária --- e dois métodos de sensibilidade
independentes: o detector sequencial de mudança de regime de Rodionov
(STARS, na sigla inglesa) \cite{Rodionov2004} e a
divergência entre matrizes de transição anuais (Kullback-Leibler e variação
total, com significância por permutação). Uma fronteira só é confirmada
quando o método primário e ao menos um método de sensibilidade convergem;
divergência entre os três seria documentada como sensibilidade ao método, e
não resolvida por escolha.

O procedimento devolveu duas fronteiras robustas --- \textbf{2001}
($F = 62{,}2$) e \textbf{2020} ($F = 21{,}5$) --- e uma terceira, em 1991,
com suporte bem mais fraco ($F = 10{,}3$), não adotada. Daí resultam os três
atos que organizam a narrativa do Capítulo~\ref{cap:resultados}: 1985--2000,
2001--2019 e 2020--2024. Dois pontos merecem registro. Primeiro, o p-valor
do sup-F é bruto, porque muitos anos-candidatos são testados; por isso a
periodização passou por uma verificação de falso positivo, aplicando o mesmo
procedimento a séries de ruído para medir quantas quebras aparecem por acaso.
Segundo, a ausência de quebra também é resultado: o Código Florestal de 2012
não produz quebra empírica em nenhuma das três séries, o que rebaixa seu
estatuto de marco periodizador --- ainda que não diga nada sobre seus
efeitos por outras vias.

A tipologia dos marcos institucionais (A, B ou C) organiza essa assimetria:
marcos com quebra empírica coincidente e mecanismo plausível, marcos com
suporte parcial e marcos de importância discursiva sem suporte empírico
detectável são tratados como categorias distintas ao longo do texto.

\section{Instrumental analítico}
\label{sec:met-instrumental}

Os métodos empregados organizam-se pelas quatro frentes de investigação
enunciadas nos objetivos específicos; a definição operacional de cada métrica
consta do glossário (Apêndice~\ref{ap:glossario}).

\subsection{As três camadas de exigência}
\label{sec:met-camadas}

Os instrumentos obedecem a uma hierarquia. Eles não respondem à mesma
pergunta, e tratá-los como equivalentes seria a origem mais provável de um
excesso de afirmação. O trabalho distingue três camadas, e cada afirmação do
Capítulo~\ref{cap:resultados} declara em qual delas se apoia
(Quadro~\ref{quadro:camadas}).

\begin{quadro}[htb]
\centering
\caption[Camadas de exigência por tipo de desenho]{Camadas de exigência: o que cada desenho consegue responder.}
\label{quadro:camadas}
\small
\begin{tabular}{p{2.6cm}p{5.0cm}p{6.0cm}}
\toprule
\textbf{Camada} & \textbf{Pergunta que responde} & \textbf{Limite} \\
\midrule
Correlação em nível estadual &
Houve coincidência temporal entre as duas séries no estado? &
Poucas observações (de 11 a 39) e nenhum controle por confundidores \\
\addlinespace
Painel com efeitos fixos &
Dentro de cada unidade, a variação de $X$ acompanha a de $Y$ depois de
removidos os choques comuns a todos no mesmo ano? &
Não isola marcos específicos; não controla determinante que varie dentro
da unidade ao longo do tempo \\
\addlinespace
Diferenças em diferenças &
Este marco regulatório teve efeito diferencial em Goiás frente a estados
de controle? &
\textbf{Não identifica neste caso}: os marcos testados são federais, de
modo que não existe grupo não tratado \\
\bottomrule
\end{tabular}
\fontefig{elaboração própria.}
\end{quadro}

A terceira camada exige explicação, porque foi \emph{rebaixada} no curso da
pesquisa e isso altera o que o trabalho pode dizer sobre política. Os quatro
marcos submetidos ao teste --- Plano Real, Lei Kandir, Código Florestal e
Cerrado Manifesto --- são federais: valeram no mesmo dia para o estado
tratado e para os controles do mesmo bioma. Sem grupo não tratado, nem o par
que sobrevive a todos os testes de robustez --- tendências paralelas, estudo
de evento e placebo temporal --- identifica efeito causal. O problema está no
desenho, não na estimação, e nenhum teste adicional o corrige; separar as
hipóteses exigiria variação \emph{dentro} do estado, entre municípios mais e
menos expostos a cada marco, que este desenho não possui. O que o coeficiente
mede é exposição diferencial ao mesmo choque, e é assim que ele é reportado:
como sensibilidade de co-movimento. Os marcos institucionais aparecem no
Capítulo~\ref{cap:resultados} como contexto da periodização, nunca como motor
de um resultado.

\subsection{Padrão espaço-temporal: centros de massa}

A descrição do deslocamento usa o \emph{centro médio} ponderado --- a média
das coordenadas de cada unidade, ponderada pela área da classe ou pelo
efetivo do rebanho ---, calculado ano a ano sobre os centroides das unidades
em coordenadas projetadas (EPSG:5880), com os pesos vindos da contagem de
pixels descrita na Seção~\ref{sec:met-area}. É um
instrumento clássico da estatística espacial descritiva, na linhagem de
\citeonline{Lefever1926}, e não uma construção \emph{ad hoc} (Seção~\ref{sec:ref-centroide}).
Duas peças o acompanham: o \emph{centro mediano}, obtido pelo algoritmo de
\citeonline{Weiszfeld1937}, que minimiza a soma das distâncias e é robusto a
um aglomerado dominante; e a elipse de desvio-padrão \cite{Yuill1971}, que
descreve dispersão e orientação da mancha.

Um centro de massa é uma estatística pontual, e sem margem de erro não há
como distinguir deslocamento real de ruído. A decisão D19 supre isso: cada
deslocamento é acompanhado de um intervalo de confiança de 95\% obtido por
\emph{bootstrap} das AMCs, com reamostragem com reposição e $B = 2000$
replicações. A regra operacional, aplicada em todo o texto: \textbf{um
deslocamento cujo intervalo inclui zero nunca é reportado em quilômetros} ---
descreve-se a variável como ancorada.

\subsection{Mecanismo: transições e painel com efeitos fixos}

O mecanismo por trás do deslocamento é lido em dois instrumentos
complementares. O primeiro é a matriz de transição calculada pixel a pixel
entre pares de anos, e não inferida da diferença entre estoques: só ela
permite dizer que classe deu origem a qual, decompondo cada hectare que
mudou em um fluxo orientado. A matriz primária opera sobre sete grupos de
classes, incluindo o Mosaico de Usos como grupo próprio, pelas razões
expostas na Seção~\ref{sec:met-decisoes}.

O segundo é o modelo de painel com efeitos fixos de dois sentidos, principal
ferramenta de inferência do trabalho (decisão D8):

\begin{equation}
\Delta \mathrm{lulc}_{it} = \alpha_i + \gamma_t + \beta \,\Delta x_{it} + \varepsilon_{it}
\label{eq:2fe}
\end{equation}

\noindent onde $\Delta \mathrm{lulc}_{it}$ é a variação anual do uso da
terra na unidade $i$; $\alpha_i$ é o efeito fixo de unidade, que absorve
tudo o que não muda no tempo naquele lugar --- solo, relevo, distância a
mercados, tradição produtiva --- mesmo que tais atributos não estejam
observados no painel; $\gamma_t$ é o efeito fixo de ano, que absorve os
choques comuns a todas as unidades (preços internacionais, câmbio, crises);
$\Delta x_{it}$ é a variação da explicativa de interesse; e $\varepsilon_{it}$
é o erro, com erros-padrão agrupados por unidade. O estimador compara cada
unidade consigo mesma ao longo do tempo, e não umas com as outras.

A especificação em diferenças não é a do manual, que costuma escrever o
modelo em níveis: é escolha deliberada, coerente com D7, porque o objeto
modelado é a \emph{mudança} do uso da terra. Nela, $\alpha_i$ absorve a
tendência própria de cada unidade e $\gamma_t$, o choque comum às variações
daquele ano.

Os diagnósticos que acompanham o painel são o fator de inflação de variância
(VIF), para multicolinearidade entre as explicativas, e o exame da estrutura
espacial dos resíduos. Este último não é rotina: a autocorrelação espacial
dos resíduos (I de Moran, com significância por permutação) mostrou-se
significativa em 115 dos 140 testes na malha municipal e em 125 dos 140 na
malha de AMC --- a dependência espacial não é, portanto, artefato da unidade
escolhida. O que o painel não explica ainda possui organização geográfica, e
foi isso que levou a dimensão espacial ao desenho da frente seguinte, com
modelos de defasagem espacial (SAR, \emph{spatial autoregressive}) e de erro
espacial (SEM, \emph{spatial error model}), além do diagnóstico local de
aglomeração (LISA, \emph{local indicators of spatial association}), este
último tratado como exploratório por não incorporar correção para testes
múltiplos.

Essa dependência tem consequência direta sobre a inferência do painel, e
convém declará-la: o erro-padrão agrupado por unidade, usado nas
especificações principais, corrige a correlação \emph{dentro} de cada
unidade ao longo do tempo, mas não a correlação \emph{entre} unidades
vizinhas no mesmo ano. O reestimador espacial foi rodado justamente para
medir o tamanho desse resíduo: os parâmetros de defasagem e de erro espacial
situam-se entre $+0{,}35$ e $+0{,}56$, confirmando a dependência, e os
coeficientes de interesse atenuam-se em no máximo 12\% quando ela é
modelada. A dependência espacial é, portanto, real e mensurada, e não
derruba os sinais reportados --- mas as magnitudes do painel devem ser lidas
como limite superior.

Nas séries agregadas no nível estadual, em que o painel não se aplica, a
inferência usa erros-padrão robustos a heterocedasticidade e autocorrelação
(HAC, na receita de \citeonline{NeweyWest1987}, com duas defasagens), porque,
com cerca de 38 observações anuais, o erro-padrão ingênuo subestima a
incerteza --- embora, nesse tamanho de amostra, o próprio teste corrigido
tenha pouco poder, o que protege contra falso positivo e esconde efeitos
reais fracos.

\subsection{Deslocamento: precedência temporal e transbordamento}

O teste da hipótese de deslocamento intra-estadual --- a lavoura do Sul
empurrando o pasto para o Norte --- começa pela correlação cruzada com
defasagens (\emph{lead-lag}),
instrumento descritivo que localiza em que encaixe temporal duas séries
melhor se ajustam e orienta onde aplicar os testes formais. O teste formal
de precedência é o de Granger, cujo nome deve ser lido com rigor: ele mede
\emph{precedência preditiva} --- se o passado de uma série ajuda a prever a
outra --- e não causalidade.

A decisão D16 vem de um erro próprio: um teste de Granger aplicado a séries
integradas fabrica precedência inexistente. O exame formal das ordens de
integração, por triangulação entre os testes de Dickey-Fuller aumentado (ADF)
e de Kwiatkowski-Phillips-Schmidt-Shin (KPSS), revelou que a série de
pastagem do Norte é integrada de ordem dois, I(2) --- nem sua primeira
diferença é estacionária --- enquanto a série de
agricultura do Sul é I(0). Séries de ordens diferentes não podem ser
cointegradas, de modo que não há relação de longo prazo entre os níveis para
o teste detectar, e qualquer ``precedência'' encontrada nessa configuração é
artefato.

Três instrumentos corrigem o problema. O primeiro é o procedimento de
Toda-Yamamoto \cite{TodaYamamoto1995}, que ajusta um modelo vetorial autorregressivo em níveis com
$p + d_{\max}$ defasagens --- $p$ defasagens ótimas mais $d_{\max}$
defasagens extras, iguais à ordem máxima de integração --- e aplica o teste
de Wald apenas às $p$ primeiras, tornando a inferência válida sob
integração. O segundo são os \emph{placebos direcionais}: o mesmo teste é
repetido contra destinos economicamente implausíveis; se o sinal persiste
onde não há mecanismo, ele não é direcionado, e sim co-tendência genérica. O
terceiro é o teste de transbordamento direcional entre vizinhos, executado
em tempo contínuo justamente para não reimportar as fronteiras da
periodização (princípio anticircularidade da decisão D12).

\subsection{Drive comum: desenho \emph{shift-share} e inferência por permutação}

A hipótese de que um mesmo motor macroeconômico atinge diferentemente
lugares diferentes é testada por um desenho de interação do tipo
\emph{shift-share} (ou Bartik): o \emph{shift} é o choque nacional --- a
taxa de câmbio ---, comum a todas as unidades em cada ano; a \emph{share} é
a exposição local, medida por uma característica física exógena, o gradiente
de aptidão edafoclimática. O desfecho da especificação confirmatória é a
\textbf{variação anual do rebanho bovino}; a mesma grade foi rodada também
sobre variação de área de pastagem e de área de agricultura, e a
Seção~\ref{sec:res-perna3} reporta os três. O interesse recai sobre o
coeficiente da interação, isto é, sobre o \emph{gradiente} do efeito, e não
sobre seu nível: com efeito fixo de ano no modelo, o nível do choque comum é
absorvido por construção. A previsão testada é que a depreciação cambial
faça o rebanho crescer \emph{mais onde a aptidão edafoclimática é baixa} ---
a fronteira do Norte --- e menos no núcleo apto do Sul, onde a mesma terra
tem uso alternativo mais rentável em lavoura; o coeficiente esperado da
interação é, portanto, negativo. Essa divisão de trabalho com a literatura é
explícita (Seção~\ref{sec:ref-cambio}).

A inferência exigiu tratamento próprio. Para desenhos \emph{shift-share} com
um único \emph{shifter}, o erro-padrão agrupado é otimista
\cite{Adao2019, Borusyak2022}. O trabalho substituiu-o por
inferência por permutação do \emph{shifter}: a série anual do câmbio é
reembaralhada entre os anos, mantidas fixas as \emph{shares}, a interação e
o coeficiente são recomputados a cada replicação, e o p-valor é a fração de
coeficientes nulos que igualam ou superam o observado em módulo. Duas
variantes foram usadas, e a distinção importa: a permutação livre quebra a
autocorrelação do câmbio e tende a ser otimista; a permutação circular, que
rotaciona a série sobre suas posições, preserva a autocorrelação do choque
macroeconômico e é a versão defensável para uma série serialmente
correlacionada. A bateria completa-se com placebos de desfecho (a mesma
interação sobre variáveis que não deveriam responder, como área urbana e
água), placebo temporal e \emph{jackknife} por ano.

O desenho tem um teto: o \emph{shifter} é uma série nacional, com cerca de 38
realizações, de modo que o poder estatístico é limitado pela dimensão
temporal e não aumenta com o número de unidades espaciais. Nenhuma quantidade
adicional de AMCs resolve isso; a consequência para a leitura dos resultados
está na Seção~\ref{sec:met-limites}.

\subsection{Teto de oferta e formas de distribuição}

A quarta frente examina a hipótese de esgotamento da terra convertível por
decomposição entre componentes de oferta e de demanda e pelo exame da
distribuição da idade da pastagem no momento de sua conversão para a
lavoura. Nesse último ponto, a pergunta é se há um ou dois mecanismos
operando, e ela é respondida por mistura de gaussianas: ajustam-se modelos
de uma e de duas componentes e comparam-se pelo critério de informação
bayesiano (BIC). Aqui, porém, entra uma restrição que o próprio desenho
impõe (decisão D23): como a análise roda sobre o \textbf{censo} de eventos, e
não sobre uma amostra, o $n$ é a população inteira, e qualquer afastamento
ínfimo do modelo de uma componente produz diferença de BIC astronômica. A
magnitude do $\Delta$BIC mede o tamanho de $n$, não a força da evidência, e
por isso ela \emph{não} é usada aqui como critério probatório. O que sustenta
a leitura de duas populações são evidências que não dependem de ajuste ---
as medianas de idade por origem anterior e o contraste entre as frações de
pasto jovem e velho ---, com o ajuste servindo apenas para descrever a forma.
É um teste de \emph{forma} da distribuição, não de nível, e as duas leituras
não são intercambiáveis --- ponto a que a Seção~\ref{sec:met-decisoes}
retorna.

\subsection{Multiplicidade}

Sempre que uma família de testes é executada --- e algumas grades deste
trabalho passam de uma centena ---, a significância bruta deixa de ser
informativa: com corte de 5\%, cerca de cinco em cem resultados aparecem
significativos por acaso. Os resultados de grades exploratórias são, por
isso, reportados sob controle da taxa de falsas descobertas
\cite{BenjaminiHochberg1995}, declarando-se explicitamente quantos testes compõem a
família --- porque o tamanho da família altera o veredito, e um resultado
que sobrevive numa grade pequena pode não sobreviver quando ela é ampliada.

\section{Réguas de robustez}
\label{sec:met-robustez}

Além dos testes internos a cada método, as afirmações centrais passam por
cinco réguas transversais, cada uma respondendo a uma objeção previsível (Quadro~\ref{quadro:reguas}).

\begin{quadro}[htb]
\centering
\caption[Réguas de robustez das afirmações centrais]{Réguas de robustez aplicadas às afirmações centrais.}
\label{quadro:reguas}
\small
\begin{tabular}{p{2.2cm}p{4.6cm}p{6.8cm}}
\toprule
\textbf{Régua} & \textbf{Pergunta} & \textbf{Procedimento} \\
\midrule
Tempo & O achado depende de onde cortamos o tempo? &
Quatro resoluções temporais (decisão D12) e teste de largura de janela móvel \\
\addlinespace
Escala & O achado depende da malha espacial? &
Município $\times$ AMC $\times$ pixel; recálculo sem malha intermediária \\
\addlinespace
Latitude & É efeito próprio da variável ou o gradiente
Sul--Norte? &
Correlação parcial controlando latitude em todo recorte transversal
(decisão D14) \\
\addlinespace
Integração & A precedência estimada é espúria? &
Ordens de integração, Toda-Yamamoto e placebos direcionais (decisão D16) \\
\addlinespace
Incerteza & Qual o intervalo de confiança? &
\emph{Bootstrap} de unidades; nada é reportado como deslocamento se o
intervalo inclui zero (decisão D19) \\
\bottomrule
\end{tabular}
\fontefig{elaboração própria.}
\end{quadro}

A régua temporal é método reutilizável, e não verificação pontual. Resumir
quatro décadas em períodos é inevitável, mas qualquer recorte é uma escolha,
sujeita a três riscos: preferir
inadvertidamente o corte que embeleza o achado; comparar períodos de
durações desiguais; e circularidade, quando se ``descobre'' nos períodos
aquilo que os definiu. A decisão D12 responde aos três com o mesmo remédio:
toda análise por período é executada em quatro resoluções --- a série
contínua ano a ano, os atos, uma grade regular de cinco anos e as décadas
--- e o que se reporta é a \emph{concordância} entre elas. As duas grades
regulares são exógenas às fronteiras dos atos, o que as torna um teste
válido; a série contínua é a verdade de referência, e um achado que não
apareça nela é suspeito de ser artefato de agrupamento. Todas as comparações
entre períodos são anualizadas, sob pena de um período curto parecer menor
apenas por ser curto.

A leitura das discordâncias é padronizada: um achado presente em todas as
réguas é robusto; um achado presente nas réguas finas que se dilui na grossa
indica fenômeno curto no tempo, e a resposta correta é localizá-lo, não
descartá-lo; um achado presente nos atos e ausente das grades exógenas é
rebaixado, por suspeita de ser artefato das fronteiras escolhidas.

A régua de escala foi verificada de forma direta. O centro de massa
calculado sobre os centroides das AMCs assume distribuição uniforme da
classe dentro de cada unidade --- suposição não necessariamente neutra,
já que as unidades do Norte são maiores e mais irregulares que as do Sul.
O mesmo centro foi então recalculado sem nenhuma malha administrativa
intermediária, com cada pixel de 30~m pesando por si. As duas medições
concordam dentro de dois quilômetros para todas as classes, contra
deslocamentos medidos em dezenas de quilômetros: o problema da unidade de
área modificável --- MAUP, na sigla inglesa de \emph{modifiable areal unit
problem}, pelo qual resultados calculados sobre unidades espaciais mudam
quando muda o recorte dessas unidades --- existe no desenho, mas não é
material para esta métrica.

\section{Decisões metodológicas críticas}
\label{sec:met-decisoes}

Esta seção registra quatro decisões que alteraram resultados já obtidos. Elas
constam aqui, e não em nota de rodapé, porque respondem à pergunta legítima
``por que seus números mudaram?'' e porque impedem que a escolha mais simples
seja restabelecida por conveniência --- inclusive pelo próprio autor, no
futuro. O registro completo está no Apêndice~\ref{ap:decisoes}.

\subsection{Censo de pixels em lugar de amostra (D21--D24)}

A idade da pastagem é medida sobre o censo completo das conversões de
pastagem em lavoura no estado --- 44.639.028 eventos ---, e não sobre
amostra: a amostragem original
sorteava pixels sobre o retângulo envolvente de Goiás, e não sobre o polígono
estadual, e confundia num único código ``dado censurado'' com ``busca no
dicionário de classes falhou'', que passaram a ser estados distintos (decisão
D22). Como o censo traz contagens, todo módulo que o consome é explicitamente
ciente de peso (contrato D24), sob pena de o gradiente estimado desaparecer
sem aviso. A troca da amostra pelo censo tem ainda uma consequência
inferencial que vale isolar (D23): com $n$ igual à população, os testes de
hipótese perdem o sentido usual, porque qualquer desvio ínfimo do modelo nulo
gera estatística arbitrariamente grande. Nos módulos que rodam sobre o censo,
portanto, a robustez é buscada na \emph{estabilidade entre recortes} --- o
mesmo padrão reaparecendo em mesorregiões, AMCs e no pixel --- e não em
p-valor ou em diferença de critério de informação (Apêndice~\ref{ap:d21}).

\subsection{A mudança de rótulo do Mosaico de Usos (D25--D26)}

No trecho final da série, parte da conversão de pastagem em agricultura passa
a ser rotulada como classe 21, ``Mosaico de Usos''. Toda medida de agricultura
em janela que alcance 2020--2024 é, por isso, reportada como o intervalo entre
a classe ``Agricultura'' isolada, que subconta, e sua união com o Mosaico, que
superconta: considera-se robusta apenas a conclusão que sobrevive nos dois
extremos, e conclusões que dependam da convenção de classe são declaradas como
tais (decisão D26). A união não é correção, e responde a uma pergunta
diferente e mais grossa --- ``quanta terra saiu de pastagem para lavoura ou uso
misto?''. Onde existe fonte independente do classificador --- área plantada,
rebanho, crédito, valor adicionado ---, é ela que responde, e o intervalo fica
reservado ao que só o sensoriamento remoto mede
(Apêndice~\ref{ap:d25}).

\subsection{Oscilação classificatória entre pastagem e savana}

O fluxo reverso de pastagem para vegetação natural, volumoso nas matrizes de
transição, é dirigido quase inteiramente à formação savânica e é
aproximadamente simétrico ao fluxo de ida --- assinatura de oscilação do
classificador, não de regeneração. Daí uma proibição de leitura adotada em
todo o texto: esse fluxo nunca é narrado como recuperação. As conclusões
centrais não dependem dele, porque se apoiam no fluxo direto de conversão e em
saldos líquidos, ambos robustos a uma oscilação simétrica, que se cancela no
líquido (Apêndice~\ref{ap:oscilacao}).

\subsection{Lotação em cabeças por hectare}

A medida de intensidade da pecuária deixou de ser expressa em Unidade
Animal e passou a ser expressa em cabeças por hectare. A conversão para
Unidade Animal aplica um fator constante ao efetivo do rebanho; sendo
constante, ela cancela em qualquer análise de tendência ou de variação, e,
sendo fixa no tempo, produz um nível não comparável entre 1985 e 2024,
porque o peso médio do animal mudou ao longo do período. Nenhuma análise
inferencial do trabalho utilizava a conversão --- ela apenas era exibida ---
e a medida reportada passou a ser a diretamente medida.

\section{Arquitetura computacional e reprodutibilidade}
\label{sec:met-repro}

O trabalho é implementado como um conjunto de 54 rotinas de coleta e análise
em Python, encadeadas por dependências explícitas, com processamento de
imagens no Google Earth Engine onde a escala exige e processamento local do
cubo de transições onde o custo de nuvem não se justifica. Cada rotina tem
documentação própria, registrando pergunta, dependências, arquivos de saída
e limitações conhecidas; o conjunto é indexado por uma ordem lógica que
descreve o papel de cada peça na construção do argumento
(Apêndice~\ref{ap:pipelines}).

Duas convenções sustentam a rastreabilidade. A primeira é que todo número
exibido --- no texto, nas figuras ou na visualização interativa --- deve ser
rastreável até um arquivo de dados que o produziu, e auditorias periódicas
percorrem os números publicados verificando essa correspondência. A segunda
é que o identificador numérico de cada rotina é permanente: rotinas não são
renumeradas quando a ordem lógica muda, porque o número é referência
cruzada em toda a documentação e apagá-lo apagaria também o registro datado
das correções.

Os produtos de uso direto previstos no mestrado profissional derivam dessa
arquitetura: o repositório de rotinas reprodutíveis sobre dados públicos e a
visualização interativa que torna os achados navegáveis.

\section{O uso de inteligência artificial como instrumento}
\label{sec:met-ia}

Esta pesquisa utilizou assistentes de inteligência artificial de forma
intensiva e sistemática, e a metodologia seria incompleta se omitisse esse
fato. A declaração é de transparência operacional; o sentido formativo dessa
escolha está na apresentação.

O escopo do uso foi a implementação: escrita e depuração de código de coleta
e análise, execução de análises exploratórias, produção de figuras e
visualizações, e auditorias de consistência entre o que os dados dizem e o
que os documentos afirmam. A formulação das perguntas de pesquisa, as
decisões metodológicas e a interpretação dos resultados permanecem do
pesquisador. A distinção não é retórica: as decisões numeradas ao longo
deste capítulo --- e as quatro descritas na
Seção~\ref{sec:met-decisoes} --- são escolhas de desenho, tomadas diante de
alternativas com custos diferentes, e é sobre elas que recai a
responsabilidade autoral.

Um instrumento dessa natureza cobra contrapartida de verificação, e a
pesquisa a organizou em quatro práticas. A primeira é o registro numerado
das decisões metodológicas, com data e justificativa, que permite reconstruir
por que cada escolha foi feita e o que ela substituiu. A segunda é a
rastreabilidade obrigatória descrita na seção anterior. A terceira é a
revisão cruzada: análises produzidas com apoio de um assistente foram
submetidas a revisão por outro, de fornecedor distinto, com verificação
manual dos números pelo pesquisador --- prática que se mostrou necessária,
porque essas revisões acertam diagnósticos estruturais e ao mesmo tempo
produzem números que parecem verificados e não são. A quarta é o registro
das autocorreções: resultados que teriam sido publicados e foram superados
por um teste adicional permanecem documentados, e não removidos.

O limite dessa prestação de contas precisa ser explícito: a verificação recai
sobre o produto --- código, número, figura --- e não sobre o processo que o
gerou. É a mesma relação que um pesquisador mantém com uma biblioteca
estatística cujo código-fonte não inspecionou: a confiança se constrói pela
validação das saídas contra fontes independentes, pela reprodutibilidade e
pela consistência interna, e não pela auditoria do instrumento. As auditorias
descritas na Seção~\ref{sec:met-decisoes} mostram o alcance e a necessidade
dessa disciplina: os defeitos ali registrados foram encontrados por
verificação sistemática --- não pelo instrumento.

\section{Limitações metodológicas}
\label{sec:met-limites}

O desenho desta pesquisa não permite algumas afirmações que seriam
convenientes, e enunciá-las aqui evita que apareçam por omissão no
Capítulo~\ref{cap:resultados}.

\textbf{Não há identificação causal forte para o papel do câmbio.} O
desenho de interação estima o gradiente do efeito, não seu nível, e a
inferência apropriada ao caso --- permutação do \emph{shifter} --- não
sustenta significância aos níveis convencionais. O resultado é corroborante e
não estabelecido. O teto é estrutural: um choque macroeconômico nacional
oferece cerca de 38 realizações independentes, e nenhuma quantidade adicional
de unidades espaciais aumenta esse número.

\textbf{Resultado nulo não é refutação.} Onde o trabalho não encontra o
canal de deslocamento intra-estadual, a afirmação sustentável é que a
assinatura testada não aparece nos dados examinados, com o poder disponível
--- e não que o fenômeno de deslocamento indireto de uso da terra tenha sido
refutado, o que exigiria desenho e escala distintos. A distinção entre ``não
encontrei'' e ``não há'' vale, em particular, quando a variável dependente é
medida com ruído, porque erro de medida no desfecho atenua qualquer
associação em direção a zero. Esta pesquisa encontrou um caso concreto disso:
um resultado nulo aparentemente limpo inverteu-se ao ser recalculado nos
mesmos municípios com desfecho medido por censo em lugar de amostra.

\textbf{A mudança de rótulo do Mosaico impõe intervalos, não pontos.} Toda
medida de agricultura em janela que alcance 2020--2024 é reportada como
intervalo entre limite inferior e superior, e conclusões que dependam da
convenção de classe são declaradas como tais. A demonstração que permitiria
converter o intervalo em correção pontual --- discriminar o artefato de
rotulagem dos sistemas integrados reais --- exigiria comparação com a coleção
anterior do MapBiomas e não foi realizada.

\textbf{Os marcos institucionais não são identificados causalmente.} Pela
razão de desenho exposta na Seção~\ref{sec:met-camadas} --- marcos federais
não deixam grupo não tratado ---, o que se estima é exposição diferencial ao
mesmo choque, e não efeito de política. Nenhuma afirmação do trabalho sobre
legislação ultrapassa esse limite.

\textbf{O terceiro ato é curto.} O período 2020--2024 tem cinco anos, e as
estatísticas calculadas sobre ele têm precisão correspondente. Onde uma
afirmação depende exclusivamente desse período, a fragilidade é declarada.

\textbf{Restrições de cobertura das fontes.} O crédito rural só existe em
série municipal a partir de 2013, de modo que os resultados que o envolvem
descrevem uma janela de cerca de oito anos, e não as quatro décadas. O IDH-M
não é publicado após 2010. O Censo Agropecuário fornece variáveis de manejo
apenas para 2017, em corte transversal. O painel com efeitos fixos controla
o que é fixo no lugar e o que é comum ao ano, mas não um determinante que
varie dentro do município ao longo do tempo e tenha ficado fora do modelo ---
razão pela qual os coeficientes estimados são descritos como associações bem
controladas, e não como efeitos causais.
